\section{The complete shared-prime obstruction for the fixed cubic character}
\label{sec:cubic-shared-local-types}

The quotient in (CER.9) can be classified at both of its
possible primes without assuming that the old character is
unramified at $269$.  Throughout this section the coefficient
prime is different from $2,3,269$, and the coefficient field
may be enlarged to contain the cubic roots of unity.

\begin{proposition}[AC-K29-local-types: the exact two bad-prime alternatives]
\label{AC-K29-local-types}
Let $f$ and $V=M_f(e)$ be as in AC-K29-bad, and let
$\rho=\operatorname{Ind}_K^{\mathbb Q}\chi$.
Thus $\det V=\det\rho=\eta_K$.
Write $C_q$ for the polynomial (CER.9), and write
$a_q(f)$ for the usual newform coefficient.
At $q=3$ one has
\[
 C_3(X)=
 \begin{cases}
  1-\dfrac{3^{2e}}{a_3(f)}X,
       &\text{if the conductor exponent of $V$ at $3$ is $1$},\\[5pt]
  1,&\text{otherwise}.
 \end{cases}                                                   \tag{CLT.1}
\]
In the first case $a_3(f)\ne0$.

At $q=269$, let $\xi$ be either nontrivial cubic character
occurring in $\rho|_{I_q}$.  Then
\[
 C_{269}(X)=
 \begin{cases}
  1-269^{2e}X^2,
       &\text{if }V|_{I_{269}}\simeq\xi\oplus\xi^{-1},\\
  1,&\text{otherwise}.
 \end{cases}                                                   \tag{CLT.2}
\]
The first case is equivalent to
$V|_{G_{\mathbb Q_{269}}}\simeq\rho|_{G_{\mathbb Q_{269}}}$;
it has tame conductor exponent two.  Conductor exponent
two by itself is not asserted to determine this inertial
type.

Consequently, for every form in AC-K29-bad's scope,
\[
 \begin{split}
 L^{\rm imp}(f,\theta_\chi,e)\ne0
 \quad\Longleftrightarrow\quad&
 \bigl(n_3(V)\ne1\ \text{or}\ a_3(f)\ne3^e\bigr)\\
 &\text{and }\quad
 V|_{I_{269}}\not\simeq\xi\oplus\xi^{-1},
 \end{split}                                                   \tag{CLT.3}
\]
where $n_3(V)$ denotes the local conductor exponent, not a
Fourier coefficient.  In particular the matching cubic
inertial type at $269$ makes the lower imprimitive value
zero, irrespective of the other local data.
\end{proposition}
\begin{proof}
We work first with the normalized representation $V$.
The cyclotomic character is unramified at the two primes
under consideration, so $M_f$ and $V$ have the same
inertial type.  Moreover
$M_f=V(-e)$, and geometric Frobenius at $q$ acts on the
Tate factor $(-e)$ by $q^e$.
We can therefore compute the extra Frobenius eigenvalues
on $(V\otimes\rho)^{I_q}/(V^{I_q}\otimes\rho^{I_q})$
and multiply each by $q^e$ to obtain $C_q(X)$.

\proofstep{The ramified quadratic prime.}
The Hensel calculation of AC-K29-ray-class gives
\[
              \rho|_{G_{\mathbb Q_3}}=1\oplus\eta_K.
                                                               \tag{CLT.4}
\]
Hence $(V\otimes\rho)^{I_3}$ can be nonzero only if
$V|_{I_3}$ contains a line of character $1$ or $\eta_K$.
Suppose it contains such a line.  Since
$\det V|_{I_3}=\eta_K$, the character of the other
one-dimensional quotient is the other one of these two
characters.
We justify that this is an actual direct sum on inertia.
Both characters are trivial on wild inertia, a pro-$3$
group.  Its image on $V$ is then unipotent and lies in a
compact subgroup of the additive coefficient field.
Such a compact subgroup is pro-$\ell$, so a continuous
image which is both pro-$3$ and pro-$\ell$ is trivial.
Thus wild inertia acts trivially.
The tame inertia quotient is abelian.  An element on
which $\eta_K=-1$ has two distinct eigenvalues, so its
eigenspaces split the full tame inertia representation.
They are the character lines $1$ and $\eta_K$.

Frobenius preserves these two lines, because it preserves
the characters $1$ and $\eta_K$ of inertia.
It follows, using the determinant, that
\[
              V=A\oplus A^{-1}\eta_K
                                                               \tag{CLT.5}
\]
with $A$ unramified.  This representation has conductor
exponent one.  Conversely, conductor exponent one means
that the inertia-fixed subspace has dimension one and the
Swan conductor is zero: the determinant is ramified, so
the fixed subspace cannot have dimension two.
Thus it gives the same decomposition by the argument
just made.
If no such line occurs, both the tensor-invariant space
and the naive invariant space vanish, and $C_3=1$.

In (CLT.5), choose a Frobenius lift on which
$\eta_K$ is one, as in the proof of AC-K29-bad.
Write $a=A(\operatorname{Fr}_3^{-1})$.
The two tensor-invariant lines in (CLT.4)--(CLT.5)
have geometric Frobenius eigenvalues $a,a^{-1}$.
The first is the naive invariant line, and the extra
line has eigenvalue $a^{-1}$.
Also $a_3(f)=3^e a$, in particular it is nonzero.
Multiplying the extra eigenvalue by $3^e$ proves
(CLT.1).

\proofstep{The inert cubic prime.}
The extension $K_{269}/\mathbb Q_{269}$ is the
unramified quadratic extension, and the cubic Kummer
character is tamely ramified of order three there.
Consequently
\[
                \rho|_{I_{269}}=\xi\oplus\xi^{-1},
                    \qquad \xi\ne\xi^{-1}.
                                                               \tag{CLT.6}
\]
A Frobenius lift exchanges these two characters, since
$269\equiv-1\pmod3$.
If $(V\otimes\rho)^{I_{269}}\ne0$, then $V$ contains
one of the character lines $\xi$ and $\xi^{-1}$.
Its Frobenius translate is a line of the other character;
the two are distinct and therefore span $V$.
This proves that $V|_{I_{269}}$ is the actual direct
sum in (CLT.6).  In particular it has finite tame
inertia and no monodromy; there are no additional wild
or unipotent directions in this two-dimensional space.
The converse follows at once by tensoring the two
character decompositions.

Choose compatible ordered character bases for $V$ and
$\rho$.  An arithmetic Frobenius lift $F$ has matrices
\[
 F|_V=\begin{pmatrix}0&b\\c&0\end{pmatrix},
 \qquad
 F|_\rho=\begin{pmatrix}0&d\\h&0\end{pmatrix}.
                                                               \tag{CLT.7}
\]
Here $bc=dh=1$, since both determinants are
$\eta_K(F)=-1$.
Thus each matrix has square one.  A diagonal change
of character basis identifies the two matrices,
proving the claimed isomorphism of local
representations.  On the two tensor-invariant lines
$V_\xi\otimes\rho_{\xi^{-1}}$ and
$V_{\xi^{-1}}\otimes\rho_\xi$, Frobenius exchanges
the lines with coefficient product $(bc)(dh)=1$.
Its characteristic polynomial is therefore $T^2-1$.
Geometric Frobenius has the same eigenvalues $1,-1$.
Since $\rho^{I_{269}}=0$, there is no naive invariant
line to subtract.  Multiplication by $269^e$ gives
(CLT.2).
In the complementary inertial types the entire invariant
space is zero, and $C_{269}=1$.
The matching type has no inertia-fixed vector and no
wild conductor, so its conductor exponent is two.

Finally AC-K29-bad proves the primitive value is
nonzero and that no other prime contributes an extra
factor.  Evaluating (CLT.1) at $3^{-e}$ gives
$1-3^e/a_3(f)$ in its first case; evaluating the first
case of (CLT.2) at $269^{-e}$ gives zero.
This proves (CLT.3).
\end{proof}

\paragraph{Scope.}
This classifies the full shared-prime factor condition of
AC-K29-lower for the fixed cubic character, including
forms ramified at $269$.  It does not exclude the matching
local type from the actual source family and does not
construct a primitive Euler class there.
The conductor notation $n_3(V)$ in (CLT.3) is deliberately
separate from the Fourier coefficient $a_3(f)$.
The argument uses the actual finite inertia characters
and their Frobenius action, not a numerical asymptotic
or a residual congruence in place of a local isomorphism.

\paragraph{Realization dictionary and source scope.}
The cubic Artin motive is the induced rank-two representation
of the finite Kummer cover.  The local calculation concerns
its tensor product with the normalized modular realization.
The discrepancy between the primitive and naive factors is
exactly the additional inertia-fixed realization in (CER.9).
No assertion about a new algebraic cycle or a comparison with
the spectral/prolate Hilbert space is involved.  The local
Kummer facts used here are proved in AC-K29-ray-class; the
primitive/naive determinant comparison and the nonzero
primitive value are the cited-source inputs of AC-K29-bad.
The new proof above supplies the complete local linear
algebra, including the nonsemisimple-inertia exclusions.
